\chapter{Data preparation and deployment tools}\label{utility_chap}
\index{Data preparation}
\index{mkdataset.py@\texttt{mkdataset.py}}
\index{mkdataset.py@\texttt{mkdataset.py}!one-hot encoding}
\index{mkdataset.py@\texttt{mkdataset.py}!reference category}
\index{mkdataset.py@\texttt{mkdataset.py}!missing values}
\index{Key file}
\index{Input structure file}
\index{neuron2web.py@\texttt{neuron2web.py}}
\index{Browser calculator}
\index{Deployment!browser calculator}
\index{Scaling factors!deployment}
\index{Label specification}
\index{Model verification!deployed calculator}
\index{Worked examples}

The current support tools are standard-library Python programs in
\texttt{tools/}.  They require Python 3 but no installed packages.
The Perl, Palm OS, and early iPhone exporters from neUROn2++ are not
part of neuron 3.0.

\begin{center}
\begin{tabular}{|>{\raggedright\arraybackslash}p{1.25in}|>{\raggedright\arraybackslash}p{2.95in}|}
\hline
\textbf{Program} & \textbf{Purpose} \\ \hline
\texttt{mkdataset.py} & Convert a delimited table into a numeric
neuron dataset plus interpretation files. \\ \hline
\texttt{neuron2web.py} & Convert a saved model and its natural-unit
interface into a self-contained browser calculator. \\ \hline
\end{tabular}
\end{center}

The maintained command reference is \path{tools/README.md}.  Both
programs are exercised against committed expected outputs by
\texttt{tests/tools/run\_tools.sh}.

\section{Preparing a dataset}

\texttt{mkdataset.py} converts an ordinary delimited file into the
numeric, outcome-last form used by the engine.  It can change
delimiters, one-hot encode categorical inputs, map a two-valued text
outcome to 0 and 1, create missing-value indicator pairs, and omit one
reference category from each categorical group.

\begin{verbatim}
python3 tools/mkdataset.py --onehot --refcat \
  --key key.txt --inputs inputs.txt \
  -o data.txt raw.csv
\end{verbatim}

The key maps encoded columns back to human variables.
\texttt{inputs.txt} groups nodes belonging to the same conceptual
variable for stepwise analysis.  The program reports the resulting
column count; the number of input nodes is one less than that count
for a single outcome.

\subsection{Current grooming behavior}

\begin{itemize}
\item The first row is treated as column labels unless
\texttt{--noheader} is supplied.
\item Python's CSV parser preserves quoted fields containing embedded
delimiters.  \texttt{--delimiter} accepts any single character,
including \texttt{\string\t} for a tab.
\item Rows with a blank outcome are dropped with a warning.
\texttt{--noelim} retains them when that behavior is intentional.
\item A numeric input containing blanks becomes a two-node indicator
pair: presence followed by value.  A blank is represented by
\texttt{0,0}; a present value by \texttt{1,value}.
\item \texttt{--onehot} with no column list detects non-numeric
columns automatically.  A comma-separated list of one-based column
numbers limits conversion to selected columns.
\item A two-level text outcome becomes one binary output.  Text input
levels become grouped indicator nodes; a blank input activates none.
\item \texttt{--refcat} omits the alphabetically first input level and
records it as the reference.  This is the appropriate form for
logistic regression and other intercept-bearing models because it
avoids a perfectly collinear full indicator group.
\item \texttt{--key} records every emitted column and its source
label.  \texttt{--inputs} records conceptual-variable groups in the
syntax accepted by stepwise regression.
\item \texttt{--quiet} suppresses the narrative report without
changing the output files.
\end{itemize}

\section{Creating a browser calculator}

\texttt{neuron2web.py} combines three artifacts:

\begin{enumerate}
\item a saved network containing architecture and weights;
\item scaling factors mapping natural values to model inputs; and
\item a label specification describing fields, units, categories, and
outcome text.
\end{enumerate}

\begin{verbatim}
python3 tools/neuron2web.py \
  --network network.txt --scales scales.txt \
  --spec spec.txt -o calculator.html
\end{verbatim}

The resulting HTML file is self-contained.  Its JavaScript, model,
scaling rules, and form are embedded, so it can be opened from disk,
emailed, or served by any static web host.  Inference remains in the
visitor's browser.

\subsection{Supported models and form fields}

The exporter currently accepts saved SimpleProp and binary logistic
models.  It refuses other network types rather than approximating
their architecture.  The label specification supports:

\begin{description}
\item[\texttt{N}] A numeric field with label and optional units.
\item[\texttt{C}] A categorical menu whose levels expand to indicator
columns.  An asterisk marks a reference level omitted by
\texttt{--refcat}.
\item[\texttt{B}] A binary choice.
\item[\texttt{K}] A missingness indicator paired with a numeric field.
\item[\texttt{O}] Labels for outcomes 0 and 1.
\item[\texttt{R}] Optional text for the reported odds.
\end{description}

The complete label grammar and validation contract are in
\path{docs/deploy.md}.  The tool checks the specification against the
saved network's input count and the scaling file before it writes the
calculator.

\subsection{Verification and preview}

Before deployment, use \texttt{--eval} with a known natural-unit row
to verify the complete scaling and forward-pass path.  A dataset that
arrived already normalized needs the original natural-to-normalized
scales; scales saved from a later training run over normalized values
would merely reproduce that normalized interface.

\texttt{--serve} writes the calculator, serves it on loopback at an
operating-system-assigned port, and opens a preview in the browser.
\texttt{--title} supplies the page title.  \texttt{--quiet}
suppresses informational output.  Neither preview nor deployment
requires an external JavaScript package or web service.

\section{Worked examples}

The illustrated Civic Choice walkthrough at
\path{docs/datasets/civic-choice/WALKTHROUGH.md} is the current
end-to-end guide.  The dataset-specific READMEs under
\path{docs/datasets/} provide further preparation and deployment
examples.
